\begin{frame}[fragile]{JavaScript and JSON}

\begin{itemize}
	\item JSON (JavaScript Object Notation) is a representation to express the idea of a dictionary in JavaScript. 
	\item JSON was invented to share information between a web server and a web browser in JavaScript. 
	\item As JavaScript becomes popular, JSON format is used outside JavaScript to make it a language-independent data format. 
\end{itemize}

\end{frame}

\begin{frame}[fragile]{AJAX}

\begin{itemize}
	\item JavaScript communicates with servers with AJAX (Asynchronous JavaScript and XML), and JSON objects represent the shared information. 
	\item Using AJAX with JavaScript is complicated, so we often use jQuery (line 2). 
	\item This is the code to request a web server with the POST function using JSON object.
\end{itemize}

\begin{Code}{}%
\begin{lstlisting}[firstnumber=1, label=glabels, xleftmargin=0pt, basicstyle=\tiny]% 

<script>
  $.ajax({
    method : 'POST',
    url : '/add',
    data : 'study'
  })
</script>
\end{lstlisting}%   
\end{Code}%

\end{frame}

\begin{frame}[fragile]{JavaScript as Prototype Language}

\begin{itemize}
	\item JavaScript is a Prototype language. 
	\item In the prototype languages, they do not instantiate a new object; instead, they copy an existing object and modify the copied object. 
	\item In this example, the new operator does not instantiate a new Abc but copies the existing Abc object and replaces the name as `Sam'.
\end{itemize}

\begin{Code}{Variables in an Object}%
\begin{lstlisting}[firstnumber=1, label=glabels, xleftmargin=0pt, basicstyle=\scriptsize]% 

function Abc() { // Abc object 
  ... 
}
var a = new Abc('Sam'); 
\end{lstlisting}%   
\end{Code}%
\end{frame}

\begin{frame}[fragile]{Function in an Object}

\begin{itemize}
	\item A JSON object can contain anything, including a function.
	\item We can use the function in an object, notice that lambda expression is used.
	\item JavaScript can mimic an OOP object with data and function this way. 
\end{itemize}

\begin{Code}{}%
\begin{lstlisting}[firstnumber=1, label=glabels, xleftmargin=0pt, basicstyle=\tiny]% 

var obj = {
  data: 'Sam',
  add: function(x,y) {return x + y;}
}
console.log(obj.add(10, 20)) 
\end{lstlisting}%   
\end{Code}%
\end{frame}

\begin{frame}[fragile]{This in an Object}

\begin{itemize}
	\item When a function has `this', it means the object that contains the function. 
	\item So, when we access any data in an object, we should prepend `this.'
	\item This feature also mimics OOP (CSC 260). 
\end{itemize}

\begin{Code}{}%
\begin{lstlisting}[firstnumber=1, label=glabels, xleftmargin=0pt, basicstyle=\scriptsize]% 

var obj = {
  data: 'Sam',
  getData: function() {return this.data;} // returns `Sam'
  t: function() {return this;}
}

console.log(obj.getData())
// t() returns this object
console.log(obj.t())
\end{lstlisting}%   
\end{Code}%

\end{frame}

\begin{frame}[fragile]{This Keyword in => Function}

\begin{itemize}
  \item We can use the arrow notation (lambda expression) in an object, but \textbf{We should be careful when we use `this' in an arrow function!}
	\item JavaScript functions do not return the object when we use `this' in the function defined with the => notation. 
	\item So, we should use the => notation only when we don't use other data in an object. 
\end{itemize}

\begin{Code}{}%
\begin{lstlisting}[firstnumber=1, label=glabels, xleftmargin=0pt, basicstyle=\tiny]% 

var obj = {
  data: 'Sam',
  t: function() {return this;},  
  t2: () => {return this;}
}
console.log(obj.t()) // returns JSON object
console.log(obj.t2()) // returns {}
\end{lstlisting}%   
\end{Code}%

\end{frame}


\begin{frame}[fragile]{Apply}

\begin{itemize}
	\item The apply high-order function applies a function to arguments. 
	\item It is one of the most frequently used in FP.
	\item As the Lisp/Scheme example shows, when we use apply function, it applies the function hi to the arguments '(1 2 3). 
	\item The result is (+ 1 2 3) == 6.  
\end{itemize}

\begin{Code}{}%
\begin{lstlisting}[firstnumber=1, label=glabels, xleftmargin=0pt, basicstyle=\scriptsize]% 

(define (hi x y z) (+ x y z))
(apply hi `(1 2 3)) ; it applies arguments (1 2 3) to the hi
\end{lstlisting}%   
\end{Code}%

\end{frame}

\begin{frame}[fragile]{}

\begin{itemize}
	\item JavaScript has the function apply() but uses a somewhat different form.
	\item Instead of the Lisp/Scheme form `apply f a', JavaScript uses `f.apply(a).'
	\item In this example, we apply the function hello() to an object person2. 
\end{itemize}

\begin{Code}{}%
\begin{lstlisting}[firstnumber=1, label=glabels, xleftmargin=0pt, basicstyle=\scriptsize]% 

var person = {
  hello : function(){
    console.log(this.name)
  }
}
var person2 = {
  name : 'Sam'
}
person.hello.apply(person2)
\end{lstlisting}%   
\end{Code}%
\end{frame}

\begin{frame}[fragile]{}

\begin{itemize}
	\item If necessary, we can give additional arguments to the function hello. 
	\item In this example, arguments `[1,2,3]' is given to the hello() function.
	\item Remember that the additional arguments should be in a list, as this idea is borrowed from Lisp/Scheme. 
\end{itemize}

\begin{Code}{}%
\begin{lstlisting}[firstnumber=1, label=glabels, xleftmargin=0pt, basicstyle=\scriptsize]% 

var person = {
  hello : function(x, y, z){
    console.log(this.name + ` Hello ${x}-${y}-${z}`)
  }
}
var person2 = {
  name : 'Sam'
}
person.hello.apply(person2, [1,2,3]);
\end{lstlisting}%   
\end{Code}%

\end{frame}

\begin{frame}[fragile]{Call}

\begin{itemize}
	\item The funcall() function in Lisp/Scheme is similar to the apply() function.
	\item The only difference is that we give each argument to the function (line 1). 
	\item JavaScript follows the same convention; when we use call() instead of apply(), we should give arguments one by one. 
\end{itemize}

\begin{Code}{}%
\begin{lstlisting}[firstnumber=1, label=glabels, xleftmargin=0pt, basicstyle=\scriptsize]% 

(funcall hi 1 2 3) ;; Lisp example
person.hello.call(person2, 1,2,3); // JavaScript
\end{lstlisting}%   
\end{Code}%

\end{frame}

\begin{frame}[fragile]{The `this' in JavaScript}

\begin{itemize}
	\item Compared to the \textbf{this} in Java that references this object, the \textbf{this} in Javascript can refer to different objects depending on where it is invoked. 
	\item This is confusing, and it can be a source of serious bugs. So, we should be careful when we use the \textbf{this} in JavaScript. 
	\item This code shows that the \textbf{this} returns different objects depending on where it is used. 
\end{itemize}

\begin{Code}{}%
\begin{lstlisting}[firstnumber=1, label=glabels, xleftmargin=0pt, basicstyle=\tiny]% 

'use strict'
console.log(this) // {Window} <- when invoked in a global context
function a() { 
  console.log(this)
} // undefined <- when invoked in a function
a() 
\end{lstlisting}%   
\end{Code}%
\end{frame}